\section{Notation and source crosswalk}
\label{sec:notation}
\label{sec:appendix}

The notation tables collect the symbols used across the certificate,
pairing, potential, and domination layers.  The crosswalk then records the
paper number and the corresponding tag in the repository companion.  Core
tags without a prefix refer to \emph{PROOF\_TSS\_DEFENDER\_ZONES.md}.
The prefix ES-P refers to \emph{ES\_POTENTIAL.md}, ES-G refers to
\emph{ES\_GLOBAL\_BOUNDARY.md}, and DOM refers to
\emph{DOMINATION.md}.  The crosswalk has no row for expository prose or an
unnumbered observation.

\subsection{Notation}

\newenvironment{appendixtable}
  {\par\medskip\noindent\begin{minipage}{\textwidth}\centering
   \captionsetup{hypcap=false}}
  {\end{minipage}\par\medskip}

\begin{appendixtable}
  \captionof{table}{Game and certificate notation.}
  \label{tab:notation-game}
  \small
  \renewcommand{\arraystretch}{1.08}
  \begin{tabular}{@{}p{.245\textwidth}p{.655\textwidth}@{}}
    \toprule
    Symbol & Meaning \\
    \midrule
    $\Cells=\mathbb Z^2$ & Axial cell lattice. \\
    $x=(q,r)$ & Cell with axial coordinates $q,r$. \\
    $d(x,y)$ & Hex distance
      $\max\{|q-q'|,|r-r'|,|(q-q')+(r-r')|\}$. \\
    $B_R(x)$, $B_R(U)$ & Closed radius-$R$ ball around a cell, and the union
      of those balls around a set. \\
    $\axisQ,\axisR,\axisS$ & The three unoriented axes represented by
      $(1,0)$, $(0,1)$, and $(1,-1)$. \\
    $W$ & Six consecutive cells on one axis. \\
    $\Omega(x)$ & The eighteen windows containing $x$. \\
    $P=(\sigma,m,b)$ & Finite position, player to place, and remaining
      placements in the current turn. \\
    $\Stones(P)$ & Occupied cells of $P$. \\
    $\Legal(P)$ & Legal empty placement cells of a nonterminal position. \\
    $\cnt_X(W,P)$ & Number of stones of player $X\in\{\A,\D\}$ in $W$. \\
    $E(W,P)$ & Empty cells of $W$ in $P$. \\
    $\ownwinnow(P)$ & Predicate that the mover can complete a window within
      the remaining placement budget. \\
    $\hit(P)$ & Family of empty sets of current opponent-threat windows. \\
    $\mhs(P)$, $\tau(F)$ & Minimum transversal size of the current threat
      family, or of a finite set family $F$. \\
    $p$ & Absolute placement index on a certificate path. \\
    $T$ & Declared absolute resolution horizon. \\
    $N$, $P_N$ & Certificate node and its labelled position. \\
    $S(N)$ & Legal defender replies represented by children of an AND node. \\
    $B(N)$ & Local upper bound on defender placements before certified
      resolution, including a LOSS remainder. \\
    $\Omega(N)$ & Live obligation roles reachable below $N$. \\
    $\Prot(N)$ & Cells carrying live attacker-placement or leaf-witness
      roles at $N$. \\
    $r_N(\rho)$, $r_N(y)$ & Remaining ordinary defender opportunities before
      role $\rho$ is checked, and the maximum rank carried by cell $y$. \\
    $E_N^{\D}(W)$ & Maximum defender exposure before attacker resolution or
      first attacker entry into $W$. \\
    \bottomrule
  \end{tabular}
\end{appendixtable}

\begin{appendixtable}
  \captionof{table}{Zone, coupling, and gate notation.}
  \label{tab:notation-zones}
  \small
  \renewcommand{\arraystretch}{1.08}
  \begin{tabular}{@{}p{.245\textwidth}p{.655\textwidth}@{}}
    \toprule
    Symbol & Meaning \\
    \midrule
    $Z_{\rm dir}(N)$ & Legal cells carrying a current direct obligation. \\
    $Z_{\rm seed}(N)$ & Legal seeds close enough to relay into a future
      protected carrier before its deadline. \\
    $Z_{\rm touch}(N)$ & Empties of touched defender-alive windows whose
      current count plus exposure can reach six. \\
    $Z_{\rm virgin}(N)$ & Legal seeds that can reach and fill an all-empty
      defender window within its exposure. \\
    $R_{\rm cert}(\mathcal C,N)$ & Mandatory certificate-relative union of
      the applicable zone terms. \\
    $R_{\rm search}(\mathcal C,N)$ & Optional search superset obtained by
      adding heuristic candidate terms. \\
    $X$, $Y$ & Real-only dismissed defender stones and ghost-only searched
      filler stones. \\
    $\widehat X$ & History of cells that have ever entered the real-only
      difference set. \\
    $K_b(N)$ & Extendable-hit kernel at remaining defender budget $b$. \\
    $\widehat B,\widehat r,\widehat E$ & Transition-inclusive D17 envelope
      labels for a selected child. \\
    $\mathcal G$ & Finite acyclic certificate DAG. \\
    $H_Q$ & Named attacker-threat windows at forcing gate $Q$. \\
    $F_Q=\{E(U,P_Q):U\in H_Q\}$ & Gate threat-empty family. \\
    $K(Q)$ & Exact-copy gate kernel
      $\{d:\tau(F_Q\setminus d)\le b-1\}$. \\
    $(Q,U,e)$ & Checkpoint role protecting empty $e$ of named gate window
      $U$ through the entry mask check. \\
    $\Prot^-(Q)$, $\Prot^+(Q)$ & Protected set immediately before and after
      checkpoint discharge at a gate. \\
    $f_N(\rho)$, $f_N(y)$ & Forced-hit-debited ordinary-opportunity clock for
      a role and the maximum such clock carried by a cell. \\
    $Q_N^{\D}(W)$ & Branch-coherent maximum of full-cost defender
      opportunities plus exact copied gate hits lying in $W$. \\
    $Z_{\rm dir}^{FH}$, $Z_{\rm seed}^{FH}$ & Direct and seed terms at an
      ordinary debited node. \\
    $Z_{\rm touch}^{FH}$, $Z_{\rm virgin}^{FH}$ & Touched and virgin
      completion terms using $Q^{\D}$. \\
    $p(Q)+b+2$ & Gate escape deadline after entry with $b$ placements left. \\
    \bottomrule
  \end{tabular}
\end{appendixtable}

\begin{appendixtable}
  \captionof{table}{Pairing, potential, and domination notation.}
  \label{tab:notation-other}
  \small
  \renewcommand{\arraystretch}{1.08}
  \begin{tabular}{@{}p{.245\textwidth}p{.655\textwidth}@{}}
    \toprule
    Symbol & Meaning \\
    \midrule
    $M$ & Partial matching used by a pairing strategy. \\
    $\Lambda=\langle(2,2),(0,6)\rangle$ & Index-twelve period lattice of the
      length-seven pairing construction. \\
    $\phi(q,r)=(q\bmod2,r-q\bmod6)$ & Quotient map with kernel $\Lambda$. \\
    $\lambda=\sqrt3$ & Base of the Hexo potential. \\
    $\Phi(P)$ & Sum of $\lambda^{-e(W)}$ over attacker-touched alive windows
      in the dynamic family. \\
    $\Psi_F(P)$ & Residual potential of a fixed finite target family $F$. \\
    $d_P(x)$ & Potential reduction, or danger, of a defender placement at
      $x$. \\
    $C_t$ & Birth mass enrolled after an attacker pair at epoch $t$. \\
    $\Phi_R$, $B_R^*$, $N_R$ & Finite-region potential, refined activation
      reserve, and coarse number of region targets. \\
    $(n_1,\ldots,n_5)$ & Counts of attacker-alive windows containing one
      through five attacker stones. \\
    $\Lambda(P)$ & Geometric legal support
      $\bigcup_{s\in\Stones(P)}B_8(s)$ in the domination layer. \\
    $\operatorname{Out}_n$ & Stopped outcome through $n$ placements after a
      compared defender reply. \\
    $b\preceq_n a$ & Defender reply $b$ is $n$-outcome-dominated by searched
      reply $a$. \\
    $G,H$ & Searched and omitted states in a causal alternating simulation. \\
    \bottomrule
  \end{tabular}
\end{appendixtable}

\clearpage
\subsection{Paper number to source tag}

The first column uses live references so the displayed paper numbers follow
the final section order.  Each table below has exactly the paper-number and
source-tag columns.

\newdimen\crosswalkremaining
\newcommand{\CrosswalkRule}{%
  \noindent\hspace*{.06\textwidth}%
  \rule{.88\textwidth}{.55pt}\par}
\newcommand{\CrosswalkRow}[2]{%
  \noindent\hspace*{.06\textwidth}%
  \begin{minipage}[t]{.58\textwidth}\raggedright\strut #1\strut\end{minipage}%
  \hspace*{.02\textwidth}%
  \begin{minipage}[t]{.28\textwidth}\raggedright\strut #2\strut\end{minipage}\par}
\newcommand{\CrosswalkContinuation}[1]{%
  \par
  \crosswalkremaining=\pagegoal
  \advance\crosswalkremaining by -\pagetotal
  \ifdim\crosswalkremaining<5\baselineskip\newpage\fi
  \noindent\begin{minipage}{\textwidth}\centering
  \textit{Table~\ref*{#1} continued}\par
  \CrosswalkRule
  \CrosswalkRow{\bfseries Paper number}{\bfseries Source tag}
  \CrosswalkRule
  \end{minipage}\par\nobreak}
\newenvironment{crosswalktable}[2]
  {\par\medskip
   \crosswalkremaining=\pagegoal
   \advance\crosswalkremaining by -\pagetotal
   \ifdim\crosswalkremaining<8\baselineskip\newpage\fi
   \begingroup\small
   \noindent\begin{minipage}{\textwidth}\centering
   \captionsetup{hypcap=false}
   \captionof{table}{#1}\label{#2}
   \CrosswalkRule
   \CrosswalkRow{\bfseries Paper number}{\bfseries Source tag}
   \CrosswalkRule
   \end{minipage}\par\nobreak}
  {\nobreak\CrosswalkRule\endgroup\medskip}

\begin{crosswalktable}{Core definitions and the geometric fact.}
  {tab:crosswalk-definitions}
  \CrosswalkRow{Definition~\ref*{def:D1}}{D1}
  \CrosswalkRow{Definition~\ref*{def:D2}}{D2}
  \CrosswalkRow{Definition~\ref*{def:D3}}{D3}
  \CrosswalkRow{Definition~\ref*{def:D4}}{D4}
  \CrosswalkRow{Definition~\ref*{def:D5}}{D5}
  \CrosswalkRow{Definition~\ref*{def:D6}}{D6}
  \CrosswalkRow{Definition~\ref*{def:D7}}{D7}
  \CrosswalkRow{Definition~\ref*{def:D8}}{D8}
  \CrosswalkRow{Definition~\ref*{def:D9}}{D9}
  \CrosswalkRow{Definition~\ref*{def:D10}}{D10}
  \CrosswalkRow{Definition~\ref*{def:D11}}{D11}
  \CrosswalkRow{Definition~\ref*{def:D12}}{D12}
  \CrosswalkRow{Definition~\ref*{def:D13}}{D13}
  \CrosswalkRow{Definition~\ref*{def:D14}}{D14}
  \CrosswalkRow{Definition~\ref*{def:D15}}{D15}
  \CrosswalkRow{Definition~\ref*{def:D16}}{D16}
  \CrosswalkRow{Definition~\ref*{def:D17}}{D17}
  \CrosswalkRow{Definition~\ref*{def:D18}}{D18}
  \CrosswalkRow{Definition~\ref*{def:D19}}{D19}
  \CrosswalkRow{Definition~\ref*{def:D20}}{D20}
  \CrosswalkRow{Definition~\ref*{def:D20a}}{D20a}
  \CrosswalkRow{Definition~\ref*{def:D21}}{D21}
  \CrosswalkRow{Fact~\ref*{fact:F1}}{F1}
\end{crosswalktable}

\begin{crosswalktable}{Core lemmas and results.}
  {tab:crosswalk-core-results}
  \CrosswalkRow{Lemma~\ref*{lem:L1}}{L1}
  \CrosswalkRow{Lemma~\ref*{lem:L2}}{L2}
  \CrosswalkRow{Lemma~\ref*{lem:L3}}{L3}
  \CrosswalkRow{Lemma~\ref*{lem:L4}}{L4}
  \CrosswalkRow{Lemma~\ref*{lem:L5}}{L5}
  \CrosswalkRow{Lemma~\ref*{lem:L6}}{L6}
  \CrosswalkRow{Lemma~\ref*{lem:L7}}{L7}
  \CrosswalkRow{Lemma~\ref*{lem:L9prime}}{L9$'$}
  \CrosswalkRow{Lemma~\ref*{lem:L10}}{L10}
  \CrosswalkRow{Lemma~\ref*{lem:L11}}{L11}
  \CrosswalkRow{Lemma~\ref*{lem:L12}}{L12}
  \CrosswalkRow{Lemma~\ref*{lem:L13}}{L13}
  \CrosswalkRow{Lemma~\ref*{lem:L14}}{L14}
  \CrosswalkRow{Lemma~\ref*{lem:L15}}{L15}
  \CrosswalkRow{Lemma~\ref*{lem:L16}}{L16}
  \CrosswalkRow{Lemma~\ref*{lem:L17}}{L17}
  \CrosswalkRow{Theorem~\ref*{thm:T1}}{T1}
  \CrosswalkRow{Theorem~\ref*{thm:T2}}{T2}
  \CrosswalkRow{Theorem~\ref*{thm:T3}}{T3}
  \CrosswalkRow{Theorem~\ref*{thm:T4}}{T4}
  \CrosswalkRow{Theorem~\ref*{thm:T5}}{T5}
  \CrosswalkRow{Theorem~\ref*{thm:T6}}{T6}
  \CrosswalkContinuation{tab:crosswalk-core-results}
  \CrosswalkRow{Theorem~\ref*{thm:T7}}{T7}
  \CrosswalkRow{Theorem~\ref*{thm:T8}}{T8}
  \CrosswalkRow{Corollary~\ref*{cor:T8.1}}{T8.1}
  \CrosswalkRow{Theorem~\ref*{thm:T8.2}}{T8.2}
  \CrosswalkRow{Theorem~\ref*{thm:T9}}{T9}
  \CrosswalkRow{Theorem~\ref*{thm:T10}}{T10}
  \CrosswalkRow{Theorem~\ref*{thm:T11}}{T11}
  \CrosswalkRow{Theorem~\ref*{thm:T11.1}}{T11.1}
\end{crosswalktable}

\begin{crosswalktable}{Potential- and domination-layer results.}
  {tab:crosswalk-potential}
  \label{tab:crosswalk-domination}
  \CrosswalkRow{Definition~\ref*{def:ES-potential}}{ES-P D3--D4}
  \CrosswalkRow{Lemma~\ref*{lem:ES-round}}{ES-P L5}
  \CrosswalkRow{Theorem~\ref*{thm:ES-fixed-family}}{ES-P T1}
  \CrosswalkRow{Theorem~\ref*{thm:ES-one-cycle}}{ES-P T2}
  \CrosswalkRow{Theorem~\ref*{thm:ES-cumulative}}{ES-P T3}
  \CrosswalkRow{Theorem~\ref*{thm:ES-finite-region}}{ES-P T4}
  \CrosswalkRow{Theorem~\ref*{thm:ES-greedy-refutation}}{ES-G T1}
  \CrosswalkRow{Proposition~\ref*{prop:ES-reachable-greedy}}{ES-G P1}
  \CrosswalkRow{Lemma~\ref*{lem:ES-clean-escape}}{ES-G L1}
  \CrosswalkRow{Theorem~\ref*{thm:ES-five-placement}}{ES-G T2}
  \CrosswalkRow{Lemma~\ref*{lem:ES-one-greedy}}{ES-G L2}
  \CrosswalkRow{Theorem~\ref*{thm:ES-three-pair}}{ES-G T3}
  \CrosswalkRow{Proposition~\ref*{prop:ES-no-static-pairing}}{ES-G P2}
  \CrosswalkRow{Proposition~\ref*{prop:ES-no-static-damping}}{ES-G P3}
  \CrosswalkRow{Theorem~\ref*{thm:ES-finite-reduction}}{ES-G T4}
  \vspace{.35\baselineskip}
  \CrosswalkRow{Lemma~\ref*{lem:Dom-L7}}{DOM L7}
  \CrosswalkRow{Proposition~\ref*{prop:Dom-P1}}{DOM P1}
  \CrosswalkRow{Proposition~\ref*{prop:Dom-P2}}{DOM P2}
  \CrosswalkRow{Proposition~\ref*{prop:Dom-P3}}{DOM P3}
\end{crosswalktable}
